% The triangle of orbit 5 with the local flag complex above it, as
% figures/triangle_projection.tex draws orbit 21: the figure of
% figures/triangle_cell_5.tex with, on its left, the incidence graph of
% figures/triangle_incidence_5.tex and a two-headed arrow onto the cell for
% the projection, a surjection. Hand-placed through the hook \gtwocelloverlay
% of julia/to_tex/triangle_cell_figure.jl; the triangle has the options and
% hence the coordinates of figures/triangle_cell.tex. The graph is on the
% left, at 0.64 of its size, the middle of its height (-0.439 in its own
% coordinates) at that of the picture of the triangle, 0.9452, and its right
% end 0.25 short of the arrow, which is 2.4 long, horizontal at that height,
% and stops 0.15 short of the left side of the triangle. The whole figure is drawn at
% 0.95.
\begingroup
\def\gtwocelloverlay{%
  \begin{scope}[shift={(-6.843,1.384)},scale=.64,transform shape]
    % The local flag complex above the triangle of orbit 5 (type \mathfrak D), as
% the incidence graph of its maximal cells, 10 tetrahedra and 1 triangular
% prism, and of the triangles that are common facets of several of them, each
% drawn as such, with its vertices. Hand-drawn; it is the tree of the cell
% with every edge subdivided and its seven leaves deleted: a triangle at every
% inner vertex, a cell at the middle of every inner edge and at the leaf of
% every pendant edge. The tree is that of figures/triangle_cell_5.tex, its
% centre v10 at the origin, its edges in the directions they have there, the
% inner ones of length 2.1 and the pendant ones 1.15. The triangles and the
% tetrahedra are those of figures/triangle_incidence.tex, drawn in the same
% way; every vertex is a letter (or the taxon) at its corner of the solid. The prism has its lateral edges BC,
% ED, HG, its end CDG towards node 10 and BEH towards node 8.
% Source: julia/studies/g2/atlas/flag/local_flag_cell_atlas.csv, rows 5 (QQ
% cell 13, rays 2,13,14), with the letters of figures/triangle_cell_5.tex,
% matched by the taxa at each vertex of the trees of the rays:
%   R_2^A,B,C = B,A,C; R_13^A,B,C = E,F,D; R_14^A,B,C = H,I,G; e_k = taxon k.
% Cells: pendant edges 1BEH, 2BEH (vertex v8), 6ADG, 7ADG (v9), 3CFG (v11),
% 4CFI, 5CFI (v12); inner edges BEH|CDG (v8-v10, the prism), ACDG (v9-v10),
% CDFG (v10-v11), CFGI (v11-v12). Triangles, the node rows: BEH (node 8), ADG
% (9), CDG (10), CFG (11), CFI (12), each the common facet of the cells of the
% edges at its vertex; no other triangle lies in two cells, and the three
% squares of the prism lie in no other cell.
% The letters are placed so that every triangle has its vertices at the same
% corners as the face of each adjacent cell that they span: top, bottom left,
% bottom right for the faces top-left-front, top-front-right and
% top-left-right of a tetrahedron, and left to right on its bottom face
% left-front-right; on an end of the prism, the top vertex is the top and the
% two others are read from left to right.
% A scope for figures/triangle_projection_5.tex, drawn inside the picture of
% figures/triangle_cell_5.tex at the place that the caller shifts it to; its
% names start with j, apart from those of the trees.
\begin{scope}[x=1.35cm,y=1.35cm]
\definecolor{g2TaxonBlue}{HTML}{1676A3}
\definecolor{g2TaxonYellow}{HTML}{E39A24}
\definecolor{g2HeartRed}{HTML}{D7443E}
\definecolor{g2RayLabel}{HTML}{000000}
\definecolor{g2CellLabel}{HTML}{000000}
\tikzset{
  g2inc edge/.style={draw=g2RayLabel,line width=.7pt},
  g2inc label/.style={text=g2RayLabel,inner sep=0pt,font={\sffamily\bfseries\fontsize{7}{7}\selectfont}},
  g2inc solid label/.style={g2inc label,font={\sffamily\bfseries\fontsize{9.5}{9.5}\selectfont}},
  % A triangle: an equilateral one, its vertices at its corners, the top one
  % first, then the left one.
  pics/g2tri/.style n args={3}{code={
    \draw[g2CellLabel,fill=g2TaxonBlue!12,line width=.6pt,line join=round]
      (0mm,6.3mm)--(-5.456mm,-3.15mm)--(5.456mm,-3.15mm)--cycle;
    \node[g2inc solid label] at (0mm,2.12mm) {#1};
    \node[g2inc solid label] at (-2.06mm,-1.47mm) {#2};
    \node[g2inc solid label] at (2.06mm,-1.47mm) {#3};
  }},
  % A tetrahedron, as in figures/triangle_incidence.tex: its two front faces,
  % lit from the left, and its vertices at its corners: the top one, the left
  % one, the front one, the right one; a pendant one has its taxon at the top,
  % drawn by the caller.
  pics/g2tet faces/.style={code={
    \path[fill=g2TaxonBlue!8] (0mm,6.27mm)--(-7.48mm,-2.86mm)--(1.1mm,-5.39mm)--cycle;
    \path[fill=g2TaxonBlue!26] (0mm,6.27mm)--(1.1mm,-5.39mm)--(7.48mm,-1.98mm)--cycle;
    \draw[g2CellLabel,line width=.6pt,line join=round] (0mm,6.27mm)--(-7.48mm,-2.86mm)--(1.1mm,-5.39mm)--(7.48mm,-1.98mm)--cycle;
  }},
  pics/g2tet/.style n args={4}{code={
    \pic {g2tet faces};
    \node[g2inc solid label] at (.22mm,3.3mm) {#1};
    \pic {g2tet base={#2}{#3}{#4}};
  }},
  pics/g2tet base/.style n args={3}{code={
    \node[g2inc solid label] at (-4.29mm,-1.98mm) {#1};
    \node[g2inc solid label] at (.99mm,-3.19mm) {#2};
    \node[g2inc solid label] at (4.51mm,-1.43mm) {#3};
  }},
  pics/g2pendant/.style n args={3}{code={
    \pic {g2tet faces};
    \pic {g2tet base={#1}{#2}{#3}};
  }},
  % A triangular prism, as large as a tetrahedron, lying on its side with an
  % edge towards us: its two front faces, the same parallelogram up to the
  % reflection in the horizontal axis, and its right end; the front edge stops
  % short of the letters at its two ends. Its vertices at its corners: those of the left
  % end, then those of the right end, each top, front, bottom; the lateral
  % edges join the two ends at the same corners.
  pics/g2prism/.style n args={6}{code={
    \path[fill=g2TaxonBlue!8] (-7.5mm,0mm)--(3.3mm,0mm)--(7.5mm,5.83mm)--(-3.3mm,5.83mm)--cycle;
    \path[fill=g2TaxonBlue!16] (-7.5mm,0mm)--(3.3mm,0mm)--(7.5mm,-5.83mm)--(-3.3mm,-5.83mm)--cycle;
    \path[fill=g2TaxonBlue!26] (3.3mm,0mm)--(7.5mm,5.83mm)--(7.5mm,-5.83mm)--cycle;
    \draw[g2CellLabel,line width=.6pt,line join=round]
      (-7.5mm,0mm)--(-3.3mm,5.83mm)--(7.5mm,5.83mm)--(7.5mm,-5.83mm)--(-3.3mm,-5.83mm)--cycle
      (-3.65mm,0mm)--(.5mm,0mm) (7.5mm,5.83mm)--(3.3mm,0mm)--(7.5mm,-5.83mm);
    \node[g2inc solid label] at (-1.95mm,4.28mm) {#1};
    \node[g2inc solid label] at (-5.1mm,0mm) {#2};
    \node[g2inc solid label] at (-1.95mm,-4.28mm) {#3};
    \node[g2inc solid label] at (3.9mm,4.28mm) {#4};
    \node[g2inc solid label] at (1.95mm,0mm) {#5};
    \node[g2inc solid label] at (3.9mm,-4.28mm) {#6};
  }},
  g2inc dark/.style={fill=g2TaxonBlue,draw=white,line width=.4pt,inner sep=0pt},
  g2inc light/.style={fill=g2TaxonYellow,draw=white,line width=.4pt,inner sep=0pt},
}
% The vertices of the G2-tree of the cell, which is not drawn.
\coordinate (jv1) at (2.47440,1.08735);
\coordinate (jv2) at (2.82994,-0.88864);
\coordinate (jv3) at (-2.48421,-1.51785);
\coordinate (jv4) at (-1.41851,-3.74527);
\coordinate (jv5) at (0.81718,-3.68673);
\coordinate (jv6) at (-2.20665,1.63956);
\coordinate (jv7) at (-0.30247,2.66486);
\coordinate (jv8) at (2.10000,0.00000);
\coordinate (jv9) at (-1.06899,1.80756);
\coordinate (jv10) at (0.00000,0.00000);
\coordinate (jv11) at (-1.33859,-1.61808);
\coordinate (jv12) at (-0.30769,-3.44763);
% The cells of the inner edges, at their middles, the prism moved 0.032 of
% its edge towards v10 to even its gaps to the two triangles.
\coordinate (je8) at (barycentric cs:jv8=.468,jv10=.532);
\coordinate (je9) at (barycentric cs:jv9=1,jv10=1);
\coordinate (je11) at (barycentric cs:jv10=1,jv11=1);
\coordinate (je12) at (barycentric cs:jv11=1,jv12=1);
% A cell and the triangles among its facets.
\foreach \a/\b in {jv1/jv8,jv2/jv8,jv6/jv9,jv7/jv9,jv3/jv11,jv4/jv12,jv5/jv12,
                   je8/jv8,je8/jv10,je9/jv9,je9/jv10,je11/jv10,je11/jv11,je12/jv11,je12/jv12}
  \draw[g2inc edge] (\a)--(\b);
\pic at (jv8) {g2tri={E}{H}{B}};
\pic at (jv9) {g2tri={A}{G}{D}};
\pic at (jv10) {g2tri={D}{G}{C}};
\pic at (jv11) {g2tri={G}{F}{C}};
\pic at (jv12) {g2tri={I}{F}{C}};
\pic at (jv1) {g2pendant={H}{E}{B}};
\pic at (jv2) {g2pendant={H}{E}{B}};
\pic at (jv3) {g2pendant={F}{G}{C}};
\pic at (jv4) {g2pendant={F}{I}{C}};
\pic at (jv5) {g2pendant={F}{I}{C}};
\pic at (jv6) {g2pendant={G}{A}{D}};
\pic at (jv7) {g2pendant={G}{A}{D}};
\pic at (je8) {g2prism={D}{G}{C}{E}{H}{B}};
\pic at (je9) {g2tet={A}{G}{D}{C}};
\pic at (je11) {g2tet={D}{F}{G}{C}};
\pic at (je12) {g2tet={G}{F}{I}{C}};
% The taxa of the pendant tetrahedra, as in the trees, at their top corners.
\node[g2inc dark,rectangle,minimum size=3.45mm] at ([xshift=0mm,yshift=2.25mm]jv1) {};
\begin{scope}[shift={([xshift=0mm,yshift=2.25mm]jv2)},scale=.675]
\path[fill=g2TaxonBlue,draw=white,line width=.4pt,line join=round] (0mm,3.0mm)--(-3.0mm,-2.2mm)--(3.0mm,-2.2mm)--cycle;
\end{scope}
\node[g2inc dark,circle,minimum size=3.75mm] at ([xshift=0mm,yshift=2.25mm]jv3) {};
\node[g2inc light,rectangle,minimum size=3.45mm] at ([xshift=0mm,yshift=2.25mm]jv4) {};
\begin{scope}[shift={([xshift=0mm,yshift=2.25mm]jv5)},scale=.675]
\path[fill=g2TaxonYellow,draw=white,line width=.4pt,line join=round] (0mm,3.0mm)--(-3.0mm,-2.2mm)--(3.0mm,-2.2mm)--cycle;
\end{scope}
\node[g2inc light,circle,minimum size=3.75mm] at ([xshift=0mm,yshift=2.25mm]jv6) {};
\begin{scope}[shift={([xshift=0mm,yshift=2.25mm]jv7)},scale=.675]
\path[fill=g2HeartRed,draw=white,line width=.4pt]
  (0mm,-2.7mm)..controls(-.45mm,-2.1mm)and(-3.0mm,.35mm)..(-3.0mm,1.45mm)
  ..controls(-3.0mm,3.2mm)and(-.9mm,3.4mm)..(0mm,1.75mm)
  ..controls(.9mm,3.4mm)and(3.0mm,3.2mm)..(3.0mm,1.45mm)
  ..controls(3.0mm,.35mm)and(.45mm,-2.1mm)..(0mm,-2.7mm)--cycle;
\end{scope}
\end{scope}

  \end{scope}
  \draw[g2RayLabel,line width=.9pt,
    -{Stealth[length=2.2mm,width=2mm]Stealth[length=2.2mm,width=2mm]}]
    (-3.659,0.9452)--(-1.259,0.9452);
}%
\scalebox{.95}{% Coarsest G2 Grassmannian cell 13: rays 2,13,14, orbit 5, with its rays clockwise from the top corner; G2-trees with lettered inner vertices, generated by julia/to_tex/triangle_cell_figure.jl
% Bounded edges of length 1.2686 in the cell and 2.0 in the rays, before scaling.
% Letters of the inner vertices of the rays, in the order of the rays:
%   A: the side {6,7} of its edges
%   B: the side {1,2} of its edges
%   C: the side {3,4,5,6,7} and {1,2,3,4,5} of its edges
%   D: the side {3,4,5,6,7} and {1,2,6,7} of its edges
%   E: the side {1,2} of its edges
%   F: the side {3,4,5} of its edges
%   G: the side {3,4,5,6,7} and {1,2,3,6,7} of its edges
%   H: the side {1,2} of its edges
%   I: the side {4,5} of its edges
\begingroup
\definecolor{g2RayLabel}{HTML}{000000}
\definecolor{g2CellLabel}{HTML}{000000}
\definecolor{g2TaxonBlue}{HTML}{1676A3}
\colorlet{g2SquareFill}{g2TaxonBlue!10}
\tikzset{
g2square side/.style={draw=g2RayLabel,line width=.8pt},
g2square circle/.style={draw=g2RayLabel,fill=white,line width=.8pt},
g2square cell/.style={fill=g2SquareFill},
g2square name/.style={text=g2RayLabel,inner sep=0pt,font=\boldmath},
g2square cell name/.style={g2square name,text=g2CellLabel},
  g2triangle letter/.style={text=g2RayLabel,inner sep=0pt,font={\sffamily\bfseries\fontsize{8}{8}\selectfont}},
}
\begin{tikzpicture}[x=1cm,y=1cm,baseline=(current bounding box.center)]
\fill[g2square cell] (0.0,2.8667)--(2.4826,-1.4333)--(-2.4826,-1.4333)--cycle;
\draw[g2square side] (0.0,2.8667)--(2.4826,-1.4333)--(-2.4826,-1.4333)--cycle;
% Ray 2
\draw[g2square circle] (0.0,2.8667) circle[radius=1.2032];
\node[g2square name,anchor=south] at (0.0,4.1499) {\textpdfrender{TextRenderingMode=FillStroke,LineWidth=.5pt}{\bfseries\boldmath \(\beta_2\)}};
\begin{scope}[shift={(0.0,2.8667)}]
\begingroup
\definecolor{g2TaxonBlue}{HTML}{1676A3}
\definecolor{g2TaxonYellow}{HTML}{E39A24}
\definecolor{g2HeartRed}{HTML}{D7443E}
\tikzset{
  g2tree pendant/.style={draw=black,line width=.6pt,line cap=round},
  g2tree internal/.style={draw=black,line width=2pt,line cap=round},
  g2tree vertex/.style={circle,fill=black,draw=white,line width=.4pt,inner sep=0pt,minimum size=3.8pt},
  g2tree dark/.style={fill=g2TaxonBlue,draw=white,line width=.55pt,text=white,inner sep=0pt},
  g2tree light/.style={fill=g2TaxonYellow,draw=white,line width=.55pt,text=white,inner sep=0pt},
  g2tree weight/.style={circle,fill=white,draw=none,text=black,inner sep=.15pt,minimum size=8.5pt},
  g2tree aligned pendant/.style={fill=white,text=black,inner sep=.7pt,sloped,above=1pt,font={\sffamily\fontsize{7.5}{8}\selectfont}},
  g2tree aligned internal/.style={fill=white,text=black,inner sep=1pt,sloped,above=1pt,font={\sffamily\fontsize{7.5}{8}\selectfont\bfseries}},
}
\coordinate (v1) at (0.66299,0.52872);
\coordinate (v2) at (0.82674,-0.18870);
\coordinate (v3) at (-0.82674,-0.18870);
\coordinate (v4) at (-0.36793,-0.76402);
\coordinate (v5) at (0.36793,-0.76402);
\coordinate (v6) at (-0.66299,0.52872);
\coordinate (v7) at (0.00000,0.84800);
\coordinate (v8) at (0.53337,-0.02306);
\coordinate (v9) at (-0.35058,0.40263);
\coordinate (v10) at (-0.18279,-0.37958);
\draw[g2tree internal] (v8)--(v10);
\draw[g2tree internal] (v9)--(v10);
\draw[g2tree pendant] (v1)--(v8);
\draw[g2tree pendant] (v2)--(v8);
\draw[g2tree pendant] (v3)--(v10);
\draw[g2tree pendant] (v4)--(v10);
\draw[g2tree pendant] (v5)--(v10);
\draw[g2tree pendant] (v6)--(v9);
\draw[g2tree pendant] (v7)--(v9);
\node[g2tree vertex] at (v8) {};
\node[g2tree vertex] at (v9) {};
\node[g2tree vertex] at (v10) {};
\node[g2tree dark,rectangle,minimum size=2.484mm,font={\sffamily\fontsize{4.8}{4.8}\selectfont\bfseries}] at (v1) {};
\begin{scope}[shift={(v2)},scale=0.54,transform shape]
\path[fill=g2TaxonBlue,draw=white,line width=.55pt,line join=round] (0mm,3.0mm)--(-3.0mm,-2.2mm)--(3.0mm,-2.2mm)--cycle;
\end{scope}
\node[g2tree dark,circle,minimum size=2.592mm,font={\sffamily\fontsize{4.8}{4.8}\selectfont\bfseries}] at (v3) {};
\node[g2tree light,rectangle,minimum size=2.484mm,font={\sffamily\fontsize{4.8}{4.8}\selectfont\bfseries}] at (v4) {};
\begin{scope}[shift={(v5)},scale=0.54,transform shape]
\path[fill=g2TaxonYellow,draw=white,line width=.55pt,line join=round] (0mm,3.0mm)--(-3.0mm,-2.2mm)--(3.0mm,-2.2mm)--cycle;
\end{scope}
\node[g2tree light,circle,minimum size=2.592mm,font={\sffamily\fontsize{4.8}{4.8}\selectfont\bfseries}] at (v6) {};
\begin{scope}[shift={(v7)},scale=0.54,transform shape]
\path[fill=g2HeartRed,draw=white,line width=.55pt]
  (0mm,-2.7mm)..controls(-.45mm,-2.1mm)and(-3.0mm,.35mm)..(-3.0mm,1.45mm)
  ..controls(-3.0mm,3.2mm)and(-.9mm,3.4mm)..(0mm,1.75mm)
  ..controls(.9mm,3.4mm)and(3.0mm,3.2mm)..(3.0mm,1.45mm)
  ..controls(3.0mm,.35mm)and(.45mm,-2.1mm)..(0mm,-2.7mm)--cycle;
\end{scope}
\node[g2triangle letter] at (0.3288,0.139) {B};
\node[g2triangle letter] at (-0.0963,0.3437) {A};
\node[g2triangle letter] at (-0.0628,-0.1304) {C};
\endgroup
\end{scope}
% Ray 13
\draw[g2square circle] (2.4826,-1.4333) circle[radius=1.2032];
\node[g2square name,anchor=north west] at (3.5939,-2.0749) {\textpdfrender{TextRenderingMode=FillStroke,LineWidth=.5pt}{\bfseries\boldmath \(\gamma_1\)}};
\begin{scope}[shift={(2.4826,-1.4333)}]
\begingroup
\definecolor{g2TaxonBlue}{HTML}{1676A3}
\definecolor{g2TaxonYellow}{HTML}{E39A24}
\definecolor{g2HeartRed}{HTML}{D7443E}
\tikzset{
  g2tree pendant/.style={draw=black,line width=.6pt,line cap=round},
  g2tree internal/.style={draw=black,line width=2pt,line cap=round},
  g2tree vertex/.style={circle,fill=black,draw=white,line width=.4pt,inner sep=0pt,minimum size=3.8pt},
  g2tree dark/.style={fill=g2TaxonBlue,draw=white,line width=.55pt,text=white,inner sep=0pt},
  g2tree light/.style={fill=g2TaxonYellow,draw=white,line width=.55pt,text=white,inner sep=0pt},
  g2tree weight/.style={circle,fill=white,draw=none,text=black,inner sep=.15pt,minimum size=8.5pt},
  g2tree aligned pendant/.style={fill=white,text=black,inner sep=.7pt,sloped,above=1pt,font={\sffamily\fontsize{7.5}{8}\selectfont}},
  g2tree aligned internal/.style={fill=white,text=black,inner sep=1pt,sloped,above=1pt,font={\sffamily\fontsize{7.5}{8}\selectfont\bfseries}},
}
\coordinate (v1) at (0.66299,0.52872);
\coordinate (v2) at (0.82674,-0.18870);
\coordinate (v3) at (-0.82674,-0.18870);
\coordinate (v4) at (-0.36793,-0.76402);
\coordinate (v5) at (0.36793,-0.76402);
\coordinate (v6) at (-0.66299,0.52872);
\coordinate (v7) at (0.00000,0.84800);
\coordinate (v8) at (0.57565,0.17549);
\coordinate (v9) at (-0.36191,-0.48082);
\coordinate (v10) at (-0.21374,0.30534);
\draw[g2tree internal] (v8)--(v10);
\draw[g2tree internal] (v9)--(v10);
\draw[g2tree pendant] (v1)--(v8);
\draw[g2tree pendant] (v2)--(v8);
\draw[g2tree pendant] (v3)--(v9);
\draw[g2tree pendant] (v4)--(v9);
\draw[g2tree pendant] (v5)--(v9);
\draw[g2tree pendant] (v6)--(v10);
\draw[g2tree pendant] (v7)--(v10);
\node[g2tree vertex] at (v8) {};
\node[g2tree vertex] at (v9) {};
\node[g2tree vertex] at (v10) {};
\node[g2tree dark,rectangle,minimum size=2.484mm,font={\sffamily\fontsize{4.8}{4.8}\selectfont\bfseries}] at (v1) {};
\begin{scope}[shift={(v2)},scale=0.54,transform shape]
\path[fill=g2TaxonBlue,draw=white,line width=.55pt,line join=round] (0mm,3.0mm)--(-3.0mm,-2.2mm)--(3.0mm,-2.2mm)--cycle;
\end{scope}
\node[g2tree dark,circle,minimum size=2.592mm,font={\sffamily\fontsize{4.8}{4.8}\selectfont\bfseries}] at (v3) {};
\node[g2tree light,rectangle,minimum size=2.484mm,font={\sffamily\fontsize{4.8}{4.8}\selectfont\bfseries}] at (v4) {};
\begin{scope}[shift={(v5)},scale=0.54,transform shape]
\path[fill=g2TaxonYellow,draw=white,line width=.55pt,line join=round] (0mm,3.0mm)--(-3.0mm,-2.2mm)--(3.0mm,-2.2mm)--cycle;
\end{scope}
\node[g2tree light,circle,minimum size=2.592mm,font={\sffamily\fontsize{4.8}{4.8}\selectfont\bfseries}] at (v6) {};
\begin{scope}[shift={(v7)},scale=0.54,transform shape]
\path[fill=g2HeartRed,draw=white,line width=.55pt]
  (0mm,-2.7mm)..controls(-.45mm,-2.1mm)and(-3.0mm,.35mm)..(-3.0mm,1.45mm)
  ..controls(-3.0mm,3.2mm)and(-.9mm,3.4mm)..(0mm,1.75mm)
  ..controls(.9mm,3.4mm)and(3.0mm,3.2mm)..(3.0mm,1.45mm)
  ..controls(3.0mm,.35mm)and(.45mm,-2.1mm)..(0mm,-2.7mm)--cycle;
\end{scope}
\node[g2triangle letter] at (0.4359,-0.0449) {E};
\node[g2triangle letter] at (-0.1338,-0.3541) {F};
\node[g2triangle letter] at (-0.4474,0.1891) {D};
\endgroup
\end{scope}
% Ray 14
\draw[g2square circle] (-2.4826,-1.4333) circle[radius=1.2032];
\node[g2square name,anchor=north east] at (-3.5939,-2.0749) {\textpdfrender{TextRenderingMode=FillStroke,LineWidth=.5pt}{\bfseries\boldmath \(\gamma_2\)}};
\begin{scope}[shift={(-2.4826,-1.4333)}]
\begingroup
\definecolor{g2TaxonBlue}{HTML}{1676A3}
\definecolor{g2TaxonYellow}{HTML}{E39A24}
\definecolor{g2HeartRed}{HTML}{D7443E}
\tikzset{
  g2tree pendant/.style={draw=black,line width=.6pt,line cap=round},
  g2tree internal/.style={draw=black,line width=2pt,line cap=round},
  g2tree vertex/.style={circle,fill=black,draw=white,line width=.4pt,inner sep=0pt,minimum size=3.8pt},
  g2tree dark/.style={fill=g2TaxonBlue,draw=white,line width=.55pt,text=white,inner sep=0pt},
  g2tree light/.style={fill=g2TaxonYellow,draw=white,line width=.55pt,text=white,inner sep=0pt},
  g2tree weight/.style={circle,fill=white,draw=none,text=black,inner sep=.15pt,minimum size=8.5pt},
  g2tree aligned pendant/.style={fill=white,text=black,inner sep=.7pt,sloped,above=1pt,font={\sffamily\fontsize{7.5}{8}\selectfont}},
  g2tree aligned internal/.style={fill=white,text=black,inner sep=1pt,sloped,above=1pt,font={\sffamily\fontsize{7.5}{8}\selectfont\bfseries}},
}
\coordinate (v1) at (0.66299,0.52872);
\coordinate (v2) at (0.82674,-0.18870);
\coordinate (v3) at (-0.82674,-0.18870);
\coordinate (v4) at (-0.36793,-0.76402);
\coordinate (v5) at (0.36793,-0.76402);
\coordinate (v6) at (-0.66299,0.52872);
\coordinate (v7) at (0.00000,0.84800);
\coordinate (v8) at (0.47055,0.25219);
\coordinate (v9) at (-0.14116,-0.51487);
\coordinate (v10) at (-0.32938,0.26267);
\draw[g2tree internal] (v8)--(v10);
\draw[g2tree internal] (v9)--(v10);
\draw[g2tree pendant] (v1)--(v8);
\draw[g2tree pendant] (v2)--(v8);
\draw[g2tree pendant] (v3)--(v10);
\draw[g2tree pendant] (v4)--(v9);
\draw[g2tree pendant] (v5)--(v9);
\draw[g2tree pendant] (v6)--(v10);
\draw[g2tree pendant] (v7)--(v10);
\node[g2tree vertex] at (v8) {};
\node[g2tree vertex] at (v9) {};
\node[g2tree vertex] at (v10) {};
\node[g2tree dark,rectangle,minimum size=2.484mm,font={\sffamily\fontsize{4.8}{4.8}\selectfont\bfseries}] at (v1) {};
\begin{scope}[shift={(v2)},scale=0.54,transform shape]
\path[fill=g2TaxonBlue,draw=white,line width=.55pt,line join=round] (0mm,3.0mm)--(-3.0mm,-2.2mm)--(3.0mm,-2.2mm)--cycle;
\end{scope}
\node[g2tree dark,circle,minimum size=2.592mm,font={\sffamily\fontsize{4.8}{4.8}\selectfont\bfseries}] at (v3) {};
\node[g2tree light,rectangle,minimum size=2.484mm,font={\sffamily\fontsize{4.8}{4.8}\selectfont\bfseries}] at (v4) {};
\begin{scope}[shift={(v5)},scale=0.54,transform shape]
\path[fill=g2TaxonYellow,draw=white,line width=.55pt,line join=round] (0mm,3.0mm)--(-3.0mm,-2.2mm)--(3.0mm,-2.2mm)--cycle;
\end{scope}
\node[g2tree light,circle,minimum size=2.592mm,font={\sffamily\fontsize{4.8}{4.8}\selectfont\bfseries}] at (v6) {};
\begin{scope}[shift={(v7)},scale=0.54,transform shape]
\path[fill=g2HeartRed,draw=white,line width=.55pt]
  (0mm,-2.7mm)..controls(-.45mm,-2.1mm)and(-3.0mm,.35mm)..(-3.0mm,1.45mm)
  ..controls(-3.0mm,3.2mm)and(-.9mm,3.4mm)..(0mm,1.75mm)
  ..controls(.9mm,3.4mm)and(3.0mm,3.2mm)..(3.0mm,1.45mm)
  ..controls(3.0mm,.35mm)and(.45mm,-2.1mm)..(0mm,-2.7mm)--cycle;
\end{scope}
\node[g2triangle letter] at (0.3565,0.0174) {H};
\node[g2triangle letter] at (0.0623,-0.3515) {I};
\node[g2triangle letter] at (-0.1131,0.0902) {G};
\endgroup
\end{scope}
% The cell
\node[g2square cell name,anchor=north] at (0.0,-1.5133) {\textpdfrender{TextRenderingMode=FillStroke,LineWidth=.5pt}{\bfseries\boldmath \oldstylenums{5}}};
\begingroup
\definecolor{g2TaxonBlue}{HTML}{1676A3}
\definecolor{g2TaxonYellow}{HTML}{E39A24}
\definecolor{g2HeartRed}{HTML}{D7443E}
\tikzset{
  g2tree pendant/.style={draw=black,line width=.6pt,line cap=round},
  g2tree internal/.style={draw=black,line width=2pt,line cap=round},
  g2tree vertex/.style={circle,fill=black,draw=white,line width=.4pt,inner sep=0pt,minimum size=3.8pt},
  g2tree dark/.style={fill=g2TaxonBlue,draw=white,line width=.55pt,text=white,inner sep=0pt},
  g2tree light/.style={fill=g2TaxonYellow,draw=white,line width=.55pt,text=white,inner sep=0pt},
  g2tree weight/.style={circle,fill=white,draw=none,text=black,inner sep=.15pt,minimum size=8.5pt},
  g2tree aligned pendant/.style={fill=white,text=black,inner sep=.7pt,sloped,above=1pt,font={\sffamily\fontsize{7.5}{8}\selectfont}},
  g2tree aligned internal/.style={fill=white,text=black,inner sep=1pt,sloped,above=1pt,font={\sffamily\fontsize{7.5}{8}\selectfont\bfseries}},
}
\coordinate (v1) at (0.82874,0.66090);
\coordinate (v2) at (1.03342,-0.23587);
\coordinate (v3) at (-1.03342,-0.23587);
\coordinate (v4) at (-0.45992,-0.95503);
\coordinate (v5) at (0.45992,-0.95503);
\coordinate (v6) at (-0.82874,0.66090);
\coordinate (v7) at (0.00000,1.06000);
\coordinate (v8) at (0.67166,0.20437);
\coordinate (v9) at (-0.28591,0.74055);
\coordinate (v10) at (0.03742,0.19483);
\coordinate (v11) at (-0.36732,-0.29358);
\coordinate (v12) at (-0.05584,-0.84616);
\draw[g2tree internal] (v8)--(v10);
\draw[g2tree internal] (v9)--(v10);
\draw[g2tree internal] (v10)--(v11);
\draw[g2tree internal] (v11)--(v12);
\draw[g2tree pendant] (v1)--(v8);
\draw[g2tree pendant] (v2)--(v8);
\draw[g2tree pendant] (v3)--(v11);
\draw[g2tree pendant] (v4)--(v12);
\draw[g2tree pendant] (v5)--(v12);
\draw[g2tree pendant] (v6)--(v9);
\draw[g2tree pendant] (v7)--(v9);
\node[g2tree vertex] at (v8) {};
\node[g2tree vertex] at (v9) {};
\node[g2tree vertex] at (v10) {};
\node[g2tree vertex] at (v11) {};
\node[g2tree vertex] at (v12) {};
\node[g2tree dark,rectangle,minimum size=3.105mm,font={\sffamily\fontsize{6}{6}\selectfont\bfseries}] at (v1) {};
\begin{scope}[shift={(v2)},scale=0.675,transform shape]
\path[fill=g2TaxonBlue,draw=white,line width=.55pt,line join=round] (0mm,3.0mm)--(-3.0mm,-2.2mm)--(3.0mm,-2.2mm)--cycle;
\end{scope}
\node[g2tree dark,circle,minimum size=3.24mm,font={\sffamily\fontsize{6}{6}\selectfont\bfseries}] at (v3) {};
\node[g2tree light,rectangle,minimum size=3.105mm,font={\sffamily\fontsize{6}{6}\selectfont\bfseries}] at (v4) {};
\begin{scope}[shift={(v5)},scale=0.675,transform shape]
\path[fill=g2TaxonYellow,draw=white,line width=.55pt,line join=round] (0mm,3.0mm)--(-3.0mm,-2.2mm)--(3.0mm,-2.2mm)--cycle;
\end{scope}
\node[g2tree light,circle,minimum size=3.24mm,font={\sffamily\fontsize{6}{6}\selectfont\bfseries}] at (v6) {};
\begin{scope}[shift={(v7)},scale=0.675,transform shape]
\path[fill=g2HeartRed,draw=white,line width=.55pt]
  (0mm,-2.7mm)..controls(-.45mm,-2.1mm)and(-3.0mm,.35mm)..(-3.0mm,1.45mm)
  ..controls(-3.0mm,3.2mm)and(-.9mm,3.4mm)..(0mm,1.75mm)
  ..controls(.9mm,3.4mm)and(3.0mm,3.2mm)..(3.0mm,1.45mm)
  ..controls(3.0mm,.35mm)and(.45mm,-2.1mm)..(0mm,-2.7mm)--cycle;
\end{scope}
\endgroup
\ifdefined\gtwocelloverlay\gtwocelloverlay\fi
\end{tikzpicture}
\endgroup
}%
\endgroup
